% GENERATED FILE. Edit canonical lessons, then run npm run generate:books.
% canonical-source-hash: fnv1a64:ade1e78ba1d6c852
% canonical-lessons: FA-C429-rah-ahan, FA-C429-sakku, FA-C429-vagon, FA-C429-kupe, FA-C429-reyl

\chapter{The railway}
\label{ch:fa-the-railway}
\hlchaptermodality{429}

\begin{chapteropening}
\textbf{By the end of this chapter:} \emph{I can talk about the railway: the railway itself, platforms, carriages, compartments and the rails.}

Complete the last lesson of chapter 429: I can talk about the railway: the railway itself, platforms, carriages, compartments and the rails.
\end{chapteropening}

\section[\texorpdfstring{\fa{راه} \fa{آهن} (râh-âhan)}{راه آهن (râh-âhan)}]{\fa{راه} \fa{آهن} (râh-âhan) — a railway (US: a railroad)}

\label{lesson:FA-C429-rah-ahan}



\begin{quote}
Before the new one: say the Persian for a pickup, then the Persian for to run over.
\end{quote}

\subsection*{You'll want to know: \fa{راه} \fa{آهن}}
\textbf{\fa{راه} \fa{آهن}} (\emph{râh-âhan}) — \textquotedblleft{}a railway (US: a railroad)\textquotedblright{}.

\begin{quote}
\fa{این} \fa{شهر} \fa{ایستگاه} \fa{راه} \fa{آهن} \fa{ندارد}.
\end{quote}

\begin{quote}
\emph{In shahr istgâh-e râh-âhan nadârad.}
\end{quote}

\begin{quote}
This town has no railway station.
\end{quote}

Literally 'iron road'. \emph{Istgâh-e râh-âhan} — the railway station.

\begin{cousinweb}
\emph{râh-âhan} is a native compound, an iron road: \emph{râh} 'road' goes back to Middle Persian \emph{rāh} (EtymDB), and \emph{âhan} 'iron' is a native word that EtymDB lists.
\end{cousinweb}

\begin{grammarlens}[title={What you've built}]
One more word for telling, explaining and describing at B1.
\end{grammarlens}

\subsection*{Your turn}
\begin{itemize}
\raggedright
  \item \emph{Say it:} \emph{râh-âhan}
  \item \emph{Say it:} \emph{râh-âhan}, once more
  \item \emph{Say it:} a sentence of your own with \emph{râh-âhan}
  \item \emph{Recall:} say the Persian for a pickup, then the Persian for to run over, then say \emph{râh-âhan} again
\end{itemize}

\begin{tcolorbox}[breakable,colback=teal!4,colframe=teal!35!black,title={Before you move on}]
What does \fa{راه} \fa{آهن} mean? (\textquotedblleft{}A railway (US: a railroad)\textquotedblright{}.) Say it once more.
\end{tcolorbox}
\section[\texorpdfstring{\fa{سکو} (sakku)}{سکو (sakku)}]{\fa{سکو} (sakku) — a platform (at a station); a raised step}

\label{lesson:FA-C429-sakku}



\begin{quote}
Before the new one: say the Persian for to run over, then the Persian for a railway.
\end{quote}

\subsection*{You'll want to know: \fa{سکو}}
\textbf{\fa{سکو}} (\emph{sakku}) — \textquotedblleft{}a platform (at a station); a raised step\textquotedblright{}.

\begin{quote}
\fa{قطار} \fa{از} \fa{سکوی} \fa{دو} \fa{حرکت} \fa{میکند}.
\end{quote}

\begin{quote}
\emph{Qatâr az sakku-ye do harekat mikonad.}
\end{quote}

\begin{quote}
The train leaves from platform two.
\end{quote}

Also a stone step or ledge in front of a house.

\begin{cousinweb}
The origin of \emph{sakku} is not settled, and EtymDB does not list it. It first meant a raised stone seat or ledge, and later a station platform.
\end{cousinweb}

\begin{grammarlens}[title={What you've built}]
One more word for telling, explaining and describing at B1.
\end{grammarlens}

\subsection*{Your turn}
\begin{itemize}
\raggedright
  \item \emph{Say it:} \emph{sakku}
  \item \emph{Say it:} \emph{sakku}, once more
  \item \emph{Say it:} a sentence of your own with \emph{sakku}
  \item \emph{Recall:} say the Persian for to run over, then the Persian for a railway, then say \emph{sakku} again
\end{itemize}

\begin{tcolorbox}[breakable,colback=teal!4,colframe=teal!35!black,title={Before you move on}]
What does \fa{سکو} mean? (\textquotedblleft{}A platform (at a station); a raised step\textquotedblright{}.) Say it once more.
\end{tcolorbox}
\section[\texorpdfstring{\fa{واگن} (vâgon)}{واگن (vâgon)}]{\fa{واگن} (vâgon) — a carriage (of a train), a railway car}

\label{lesson:FA-C429-vagon}



\begin{quote}
Before the new one: say the Persian for a railway, then the Persian for a platform.
\end{quote}

\subsection*{You'll want to know: \fa{واگن}}
\textbf{\fa{واگن}} (\emph{vâgon}) — \textquotedblleft{}a carriage (of a train), a railway car\textquotedblright{}.

\begin{quote}
\fa{واگن} \fa{ما} \fa{آخر} \fa{قطار} \fa{بود}.
\end{quote}

\begin{quote}
\emph{Vâgon-e mâ âkhar-e qatâr bud.}
\end{quote}

\begin{quote}
Our carriage was at the end of the train.
\end{quote}

\emph{Vâgon-e ghazâkhori} — the dining car.

\begin{cousinweb}
\emph{vâgon} was borrowed from a European language: it is usually traced to French \emph{wagon}, which itself came from English \emph{wagon}. EtymDB does not list the Persian word.
\end{cousinweb}

\begin{grammarlens}[title={What you've built}]
One more word for telling, explaining and describing at B1.
\end{grammarlens}

\subsection*{Your turn}
\begin{itemize}
\raggedright
  \item \emph{Say it:} \emph{vâgon}
  \item \emph{Say it:} \emph{vâgon}, once more
  \item \emph{Say it:} a sentence of your own with \emph{vâgon}
  \item \emph{Recall:} say the Persian for a railway, then the Persian for a platform, then say \emph{vâgon} again
\end{itemize}

\begin{tcolorbox}[breakable,colback=teal!4,colframe=teal!35!black,title={Before you move on}]
What does \fa{واگن} mean? (\textquotedblleft{}A carriage (of a train), a railway car\textquotedblright{}.) Say it once more.
\end{tcolorbox}
\section[\texorpdfstring{\fa{کوپه} (kupe)}{کوپه (kupe)}]{\fa{کوپه} (kupe) — a compartment (in a train)}

\label{lesson:FA-C429-kupe}



\begin{quote}
Before the new one: say the Persian for a platform, then the Persian for a carriage.
\end{quote}

\subsection*{You'll want to know: \fa{کوپه}}
\textbf{\fa{کوپه}} (\emph{kupe}) — \textquotedblleft{}a compartment (in a train)\textquotedblright{}.

\begin{quote}
\fa{در} \fa{کوپه} \fa{ما} \fa{چهار} \fa{تخت} \fa{بود}.
\end{quote}

\begin{quote}
\emph{Dar kupe-ye mâ chahâr takht bud.}
\end{quote}

\begin{quote}
There were four berths in our compartment.
\end{quote}

Night trains in Iran have four- or six-berth \emph{kupe}.

\begin{cousinweb}
\emph{kupe} was borrowed from French and is usually traced to French \emph{coupé}, first a closed carriage. EtymDB does not list the Persian word.
\end{cousinweb}

\begin{grammarlens}[title={What you've built}]
One more word for telling, explaining and describing at B1.
\end{grammarlens}

\subsection*{Your turn}
\begin{itemize}
\raggedright
  \item \emph{Say it:} \emph{kupe}
  \item \emph{Say it:} \emph{kupe}, once more
  \item \emph{Say it:} a sentence of your own with \emph{kupe}
  \item \emph{Recall:} say the Persian for a platform, then the Persian for a carriage, then say \emph{kupe} again
\end{itemize}

\begin{tcolorbox}[breakable,colback=teal!4,colframe=teal!35!black,title={Before you move on}]
What does \fa{کوپه} mean? (\textquotedblleft{}A compartment (in a train)\textquotedblright{}.) Say it once more.
\end{tcolorbox}
\section[\texorpdfstring{\fa{ریل} (reyl)}{ریل (reyl)}]{\fa{ریل} (reyl) — a rail (of a railway); the rails}

\label{lesson:FA-C429-reyl}



\begin{quote}
Before the new one: say the Persian for a carriage, then the Persian for a compartment.
\end{quote}

\subsection*{You'll want to know: \fa{ریل}}
\textbf{\fa{ریل}} (\emph{reyl}) — \textquotedblleft{}a rail (of a railway); the rails\textquotedblright{}.

\begin{quote}
\fa{قطار} \fa{از} \fa{ریل} \fa{خارج} \fa{شد}.
\end{quote}

\begin{quote}
\emph{Qatâr az reyl khârej shod.}
\end{quote}

\begin{quote}
The train came off the rails.
\end{quote}

\emph{Az reyl khârej shodan} — to come off the rails.

\begin{cousinweb}
\emph{reyl} was borrowed from English and is usually traced to English \emph{rail}. EtymDB does not list the Persian word.
\end{cousinweb}

\begin{grammarlens}[title={What you've built}]
One more word for telling, explaining and describing at B1.
\end{grammarlens}

\subsection*{Your turn}
\begin{itemize}
\raggedright
  \item \emph{Say it:} \emph{reyl}
  \item \emph{Say it:} \emph{reyl}, once more
  \item \emph{Say it:} a sentence of your own with \emph{reyl}
  \item \emph{Recall:} say the Persian for a carriage, then the Persian for a compartment, then say \emph{reyl} again
\end{itemize}

\begin{tcolorbox}[breakable,colback=teal!4,colframe=teal!35!black,title={Before you move on}]
What does \fa{ریل} mean? (\textquotedblleft{}A rail (of a railway); the rails\textquotedblright{}.) Say it once more.
\end{tcolorbox}
